\sgaStableTarget{sga6.e14.eq.4-5}@@M0@@
where in the first member @@M1@@ is regarded as a coherent sheaf on @@M2@@, thus defining an element of @@M3@@, and in the second member @@M4@@ denotes the class of @@M5@@ in @@M6@@. Since @@M7@@ does not change on removing from @@M8@@ a closed subset of codimension @@M9@@, we are reduced to the case in which @@M10@@ is regular, hence smooth over the base field @@M11@@ when the latter is perfect. But in this latter case, formula (4.5) is precisely the Riemann--Roch formula for the inclusion @@M12@@ and the element @@M13@@ of @@M14@@, written using the theory of Chern classes with values in the Chow ring @@M15@@. This formula is established in [2] under the conditions considered here, by a method very close to the one used in this seminar for the analogous formula using @@M16@@. This therefore establishes, at least for perfect @@M17@@, formula (4.4) and consequently the fact that @@M18@@ is an isomorphism mod. torsion (which in turn implies, a posteriori, that the Riemann--Roch formula ``with denominators'' for a proper morphism of smooth algebraic schemes over @@M19@@ is essentially the same whether written using @@M20@@ as in [2] or using @@M21@@ as in this seminar!). The case with no hypothesis on @@M22@@ reduces easily to the case in which @@M23@@ is algebraically closed, using the fact that, for a finite extension @@M24@@ of degree @@M25@@ of @@M26@@, the kernel of the natural map @@M27@@ is annihilated by @@M28@@ (as is seen by using the trace homomorphism @@M29@@, satisfying @@M30@@).\SGAAuthorFootText{*}{It is therefore established, thanks to P. JOUANOLOU, for @@M31@@ smooth and quasi-projective over a field (cf. note p.5).}

\sgaStableTarget{sga6.e14.para.4-3}\textbf{4.3.} It seems quite plausible that formulas (4.4) hold as they stand and not only modulo torsion, or, equivalently, that formula 4.5. holds not only modulo torsion. This question is a very special case of the validity in @@M32@@ of the ``Riemann--Roch formula without denominators'' (3.1) for an immersion @@M33@@ (namely, the case @@M34@@). The validity of this special case may be regarded as quite plausible. Indeed, we already pointed out in no.~3 that, when @@M35@@ is of characteristic zero, the Riemann--Roch formula without denominators holds at least in @@M36@@; this implies in that case that both members of (4.5) have the same image in @@M37@@, or equivalently that the second of formulas (4.4) holds; this result will extend of itself to arbitrary characteristic once
